\chapter{Anecdotes on the Journey}

As I've encountered different problems related to computer systems in my life, I have noticed that sometimes other people don't react the way I might expect.  Sometimes people are supportive of finding and treating vulnerabilities.  Other times, people have told me directly that they didn't want any problems mentioned for fear that admitting there was a problem would make their system look weak.  Still other people have sometimes reacted in ways that were so self-centered that their reactions seemed foolish, almost comical.  Having seen some of these reactions, I noticed that our values and judgements about actors and victimization play a big role in how we might react to finding out we are the victims of attacks on computer systems.  I would like for us to think a little about how we might choose to value ourselves, the people working around us, and others affected by malfeasance on the web.  To that end, I would like to tell you a few small tales about my life and computers.

\section{Computers and Insects}

When I was a small boy, we lived in New Orleans.  In my part of the city, there was a large, open drainage canal.  This canal had steep sides covered with grass and weeds.  Most of the time the water was low in the canal.  Brown water snaked through the basin of the canal.  Sometimes there would be a muddy shopping cart in the weeds.  One time I saw an old man fishing for a turtle in the canal.  He used a big hook and a piece of meat.  He caught one large snapping turtle.  With a taut line, he pulled the turtle from the waters of the canal.  The snapping turtle was dripping with tan mud.  Its shell seemed as large as the man's torso.  He lifted it up and dropped it in a burlap sack.  He quickly gathered the top corners and held the bag shut with his fist.

I had been picking blackberries in the brambles of the canal.  As I put my blackberries in my plastic margarine tub, I asked the man why he was fishing for the turtle.  He told me he was going to eat the turtle.  He would make turtle soup.

I often saw rat traps tied with strings to trees in the canal.  Some people had houses along one side of the canal.  Along the edge of their backyards, they tied rat traps to tree trunks or metal fence posts.  These rat traps were like mouse traps, but much larger. Nobody ever bothered to put bait in them.  I lived in an apartment complex on the other side.  Every day I would walk to and from school with other kids from my neighborhood.  Every day, we would pass that canal.

One day, near 1979, a computer programmer came to my school.  Our teacher gathered us in the library.  We sat in a circle on the floor, cross-legged, "Indian style."  The programmer came in, pushing his cart.  Taking up the whole top of the cart was the computer; it was a big gray box with a keyboard and a built-in monitor.  Underneath, on the shelf, was another set of boxes.  One of them was an eight-inch floppy disk drive.

The programmer told us about computers.  Looking back, I realize that I forgot everything he said. He had shown us about files and the parts of the computer.  Then, there, on the computer screen, I saw something fascinating:  my first look at a video game.  It was Pac-Man in black-and-white.  There were no colors on the screen.

One day I was walking home from school with the other kids.  Because I was little, I had to walk to and from school with one of the big kids, Jeanie.  She would watch out for me and make sure I went the right way.

As we passed the canal every day, there was a little draw, like a V or little valley, going from the curb of the street down to where there was a big pipe.  The canal went under the road in this a big concrete pipe.  Some of the big kids would go down to the rim of the pipe and stand on its edge.  I was too little for that yet.  Instead, like most of the kids, I would run and jump across the little V at the top.  It was only a few steps, a few feet, but it felt like a big adventure when I made it across.

I was having fun and feeling great in the sun, fresh out of school, I decided to run.  I ran, and I jumped.  I tried to get across the little V, but I didn't make it.  I tumbled and fell into the dirt on the wall of the little valley that went down to the big pipe.

I jumped. I stumbled, and I fell.  I landed in a big ant pile.  There, on the far side of the V was a big ant pile.  It was about three feet around.  The puffy grains of cultivated dirt stood knee high to a man.  The ant pile was filled with hundreds of ants.  Before I knew it, before I felt myself really fall, I was swarmed and covered with ants!

They were biting me! I forgot about the big canal.  I scrambled up the dirt.  They were all over me!  Hundreds of ants!  Ants on my back. Ants on my legs.  Ants on my arms, in my hair and on my head.  I did what any kid would do:  I screamed and cried and ran.  I ran as fast as I could.  I ran across the grass into the apartment complex.  I ran down the streets to my home.  I ran as fast as I could.

I cried for my Mutti and told her about the ants.  She put me in a cold bath.  A little while later Jeanie came by with my book bag.  I had dropped it when I ran home.  She had rescued it. She told my Mutti about the ants.  I still remember my mom saying that she knew.  She had been glad Jeanie had checked up on me when I got home.

It gets hot in New Orleans.  I still remember having to wear a white, long-sleeved shirt to school.  It was so hot and humid.  I felt sweaty, having to wear the long-sleeved dress shirt.  Often in New Orleans, I would play in shorts and a short-sleeved shirt.  I would often run barefoot and get a suntan.  But, thanks to the ants, I had to go to school in a long-sleeved shirt.  I was so covered with ant bites that it looked like I had the chicken pox.  I had to wear the long-sleeved shirt to keep from scratching.

The next time a computer came into my life was some years later.  Now, my friends, and later I, had an Atari 2600.  Today I would look upon this as a computer, but back then, to me, it was "a video game."  Also, I had seen some other computers.  I had an old friend whose family was more rich; they had an IBM PCjr.  I saw Apple IIe's at school.  I had even attended a week-long summer camp at junior high.  I remember typing in my programs in BASIC.  I somehow had a book that I had gotten about computer programs.  It was about typing in programs to make a computer game.  I don't remember the title.  I don't have the book anymore.  I don't know what language it was written in.  However, the book was about adventures.  That's what I wanted:  computers and games and adventures.  I never did get many of those programs to work.

I also remember my friend next to me, Eddie, typing in a program in another way.  I asked Eddie what it was.  He said it was C.  C!  The teacher said C was advanced.  He didn't want the rest of us to be trying to type programs in C, if we didn't know what we were doing.  It was okay for Eddie to try, because he was advanced.  The next week, there was another week of computer camp.  I didn't go because I thought it was for the advanced kids.

The computers we used at school and saw around us:  they were expensive computers.  They cost thousands of dollars.  It was clear to me that I shouldn't even ask for one because there was no way we'd be able to buy one of those.  Yet, somehow, I ran across an ad in a magazine for a Commodore 64.  It was \$200.

Now, that was a lot of money back then, too.  It was a bare-bones system.  I would still need some more components.  But, it was \$200.

A friend of mine, Donald, had proudly told me at school how he had a job and made money mowing lawns and raking leaves.  That sounded pretty good to me:  making money.  I was mildly interested in lawn care.  Most of all, I knew I could do it.  I could push the lawnmower.  My dad made me mow the lawn.  I could earn some money.  I canvassed the neighborhood.  Some people said no.  Some people said they mowed their own lawn.  I didn't want to go too far up the street into Donald's territory.  I found one customer:  the man across the street.

I mowed the lawn for \$20 a pop.  I was the worst employee you ever saw.  Frequently, I would just quit working.  It would be hot, and the Zoysia grass would be thick and wet.  It would clog up my lawnmower because I would wait too long to go back and mow again.  I would cut one rectangular patch in the thick part of the front yard.  It would get hot, and I would leave to cool off.  I would come back eventually and start mowing again.

One day, I was bopping along, mowing real great, dreaming of imaginary greatness.  I had a good rhythm going. I was getting close to finishing at the end.  I bumped up against that big old Magnolia tree Mister Warner had in his yard.  Somewhere right around there, I hit a yellow jacket nest.

I don't know if you've ever hit a yellow jacket nest with a lawnmower, but I want to tell you:  those yellow jackets were angry.  They do not want to be disturbed.  All they want to do is eat fruit, rotten insects, and sleep in their enormous, maze-like hell of an underground yellow jacket nest.

I remember that I felt them buzzing around me and swarming.  I felt some heat on my arm, and looked down and saw one stinging.  I can still remember seeing one stuck in the skin of my arm, trying to pull away, as it stung me.

This time, I was older.  I was more mature.  I was experienced.  When it came to stings, I was advanced.  I didn't just cry when I ran home.  I ran home aggressively.

I started getting undressed outside.  I remember I stripped off my pants in the kitchen.  Insects were still buzzing in the legs of my jeans.  I swatted at them. I ran and got in a cold tub.  I still remember brushing at my arms and legs in the cold water.  I felt a few stingers come out.  The stings burned and throbbed and grew red.  I had a little over \$200.  I had been mowing all summer.

As I scrounged over my hoarded money, it turned out that I had just enough to get the computer and a TV set.  I remember feeling stunned as my mother asked me if I had enough for tax.  Tax! What tax, I said.  Sales tax, they said.  How much was that, I asked in a daze.  Somehow, they must have saved me, because I came home with a computer and a black-and-white TV set to use as a monitor.

It was a Commodore 64.  I would type the programs in.   It came with a manual.  I had learned a little BASIC in school.  Trouble was, every time I would turn the computer off, it would forget the programs.  You see, \gls{ram} is a volatile form of memory.  Without electricity constantly charging that type of memory, it would forget.  There was ROM in the computer, and in cartridges for games, but what was needed was a more persistent local storage.  I needed a tape drive.

I had to go back and mow some more lawns.  I did go back.  I did mow some more lawns. Mister Warner did pay me, even though I was the worst employee ever.  He did help me to realize that there were jobs out there.  Yet, once I got that tape drive, I don't think I ever went back.  Soon we moved away, and I was on to the next adventure:  high school.

After high school, I got my first chance at a real job:  joining the Army.  Now, at the time, I was only seventeen.  I remember telling my mom and dad, The recruiters said that since I was under age, you would have to sign for me.  I thought that would wrap up that effort.  Instead, my parents said, Where do I sign?  Next thing I knew, I was here and away to Basic Training.

I learned how to work in Basic Training.  I put forth effort.  I felt like I was under a lot of stress, with no way to escape, for the first time in my life.  I learned how to work and get along with people I might not otherwise be around.  After Basic, I went on to radio school.  I became a RATT operator, radio automatic teletype.  I went to jump school.  I spent my enlistment jumping out of planes and operating radios.  It was a great job for a young man to have.

This job was where I really learned about computers and telecommunications.  The Army made us do everything to send a message.  We began by learning to type on a 1949 teletype.  We learned how to generate our own electricity, and we learned how to hookup systems to commercial power if it was available.  We were taught to construct our own antennas, raise our own 50-foot antenna towers.  We used one-time pads, hardware and software \gls{encrypt}ion techniques.  We used punched paper tape, Bell Labs modems with oscilloscopes, teleprinters and shortwave radios.  We used \gls{encrypt}ed, wireless networking systems back before the Internet became popular among the public.

One day I was on the job during Operation Desert Shield in Saudi Arabia.  I had had a long day.  I had been up for most of the day.  I had worked the night shift.  I was about to go to bed when Sergeant D. told me, You have latrine detail.  Back then, we used burnout latrines.  So, what this means was that we would drag the cans of sewage out of the latrine hut, and burn off the sewage that had pooled inside.  To do this, you would use a five gallon can of diesel fuel; pour that into the can.  Then you would pour some unleaded gasoline into the can.  It would look pink over the yellowish diesel.  Immediately, a thick curtain of vapors would rise from the can.  With the cans loaded with fuel, we would get back some feet, crouch down, and strike and throw a match towards the curtain of gasoline vapor.  Whump!  The cans would catch flame.  One by one, we would set them on fire this way.

Burning sewage takes hours.  What happens is that the burning unleaded gasoline ("MOGAS," we called it in the Army) will begin to heat up and boil the diesel.  Once the diesel is hot enough, it will continue to boil and burn, too.  Diesel, by itself, will not readily catch fire at just the touch of a match.  If it's not prepared, it can actually snuff out a match.  In contrast, gasoline, of course, will immediately flash over into fire.  One of the tricks I learned on the job in the Army was that, if you could not tell the difference between MOGAS or diesel, one way to check was to dip a Styrofoam cup in the liquid.  Gasoline will eat away quickly at the Styrofoam.  It's a quick and cheap assay for telling the fuels apart.  Sometimes, in the bright sunlight, with your sense of smell dulled from being around fuel, you might not see the difference readily in the liquids themselves.  Fortunately, the Army required the cans to be kept well marked.

When you burn sewage for hours, you have to stir it.  Human waste, like the body it comes from, is loaded with water.  If you don't stir up the sewage in the can, then it won't all burn down to ash.  Since that's the goal, to neutralize the pathogens by reducing the sewage to ash, you have to stir the flaming can of burning sewage.  It smells horrible.  The foul odor permeates your hair, your clothing, your skin.  It gets in your nostrils.  It's all over you, and there's no escaping it.

By the time I was done with all this burning sewage in the 120F desert, it was time for me to break down and take a shower.  Our shower was an Australian shower, which was a canvas bag with a shower head built in.  You would get a five gallon can of water, fill the bag, hoist it up on a pole, and then shower off with the cold water.  With my bare feet in the sand, I rubbed some soap all over myself.  I rinsed off before I ran out of my five gallons.  I put on some clean clothes and walked around the camp.

I sat on my cot and made fists with my toes in the sand.  I was about to eat an MRE and get some sleep before my next shift began that evening.  I walked around the corner of the tent to get a water bottle.  I picked up an empty cardboard box to throw it away and get myself a fresh bottle from the next box.  I heard a tick and felt a bump and a sting against my big toe.  I had been stung by a yellow scorpion.

The night of the scorpion was a great adventure for me.  My friends sprang into action.  They picked me up by my clothes and carried me toward the medics.  Just then, a Blazer slid to a halt, throwing an arc of sand.  We had been asking for a support vehicle all day.  A friend hot-footed one up the hill to our spot just in time.  I remember laying on the stretcher at Battalion Aid.  Already convulsions had caused my leg to shudder rapidly.  Scorpions kill their prey with acetylcholinesterase inhibitors; those chemicals poison nerves into an electrical storm that causes electrical confusion; this leads to paralysis.  In my case, I had some kind of allergic reaction to the poison.  Instead of growing cold and immobile, my muscles twitched and convulsed.

The doctor at battalion said he only had two 500mL vials of Benedryl.  It was just enough to save one man from a death like this.  He said he needed to save his medicine for just the right moment, because he might not be able to get more.  At one point the doctor said, If the convulsions reach his chest, he'll smother.  I remember my friends all leaning in over the stretcher to check on how I was doing.  I made sure to announce I was still breathing just fine.

I didn't get any medicine.  That night, I spent one of the most painful nights of my life.  Until years later, when I was tasered, I never felt any pain even remotely close to the stabbing brush of pain inside my leg that night.  With every heartbeat, and between every heartbeat, I could imagine a tree of sharp, icy fire that was my nerves.  One of my friends took my shift as I tried to rest.

The yellow scorpion was the fourth most poisonous animal on the Arabian Peninsula.  The other three were snakes:  the carpet viper, the cobra, and the hog-nosed viper.  I never saw a carpet viper.  The cobra makes its home in irrigated areas; but, at the time, Saudi Arabia was the world's fifth leading producer of cereals.  It was heavily irrigated in some parts.  I did see a sleeping hog nosed viper, curled in a figure eight about the size of a boot print.  It was sleepy from the cold of some fog rain in the winter.  The scorpion, though, was more than poisonous enough for me.  The day after I was stung, my buddy carried me on his back the half mile to battalion aid.  I asked for some pain killers.  The medic sergeant told me they don't treat pain.  I asked for an ibuprofen.  I was denied.  We had to hobble back to camp and wait it out.  The rest of the pain subsided later that day.

Later on in those war campaigns, during Desert Storm, some people that I worked with died.  I knew one of them in passing.  He would come by our rig looking to send a message.  He was a Captain, a West Point honor graduate.  He was Puerto Rican.  He smiled a lot.  We were in a vocation where optimism wasn't openly valued.  He was a good person.  He died leading his troops from the front.

When these people died, there was nothing I could do to help them.  I had had a pretty good tour up to that point.  I was fairly good at my job.  But, when the time came, there wasn't anything I could do for these people who were so badly and suddenly wounded.  I felt terrible about this for a long time.  The day of their death was the worst day I ever had on the job.

Years later, I was in college studying Literature and Writing.  I was reading a book by Tim O'Brien called The Things They Carried.  It was about and based on O'Brien's experiences during the Vietnam War.  Somewhere in there, I remember reading a lesson about how some people believed the folly that others had died because they were stupid.  In my later tours in Iraq, I saw this mythological phenomenon rise again:  if you had an aggressive posture, the sayings went, then the enemy wouldn't \gls{attack} you.  I never found any of this to be true.

The fact is, I believe:  Bad things happen to good people who do the right thing, every day.  I know that's hard for some to accept.  Decades later, my pastor told me that I had a harsh view of theology.  Still, though it has been very difficult for me to face the world this way and accept it for what it is, I believe these thoughts.  I wish we lived in a world where justice would reign.  I wish we could see that the truth would exonerate many living souls.  Still, I see evil shadows and thieves all around.  And no, I do not see the Good always prosper.

I accept that bad things happen to good people who do the right thing, every day.  I mention all this because I want you, too, to consider it.  Sometimes I meet rationalists who just cannot accept that bad news comes to good people.  I want you to know and recognize that bad times can befall good people.

Accepting this idea can challenge some of us.  We might not want to acknowledge what this kind of inherent unfairness in life can say to our notions of justice.  It can push against our concepts of actor and victim responsibility.  Some might defensively whine that, We are all victims now.  Maybe we'll invite a straw man over so that we can knock him down.  Yet, in a world filled with environmental hazards, we have to recognize that some of those hazards carry catastrophic consequences for the wrong kind of interaction.  In line with our sense of justice or not, deep, extreme, and powerful unfairness can override the will and reckoning of our best and strongest people.  Still, many more will suffer faster and harder a more damaging fate solely because we may be weaker, slower, or somehow wrongly suited to counteract whatever overpowering \gls{threat} awaits.  Living in a world with catastrophic risks changes us all.  While some people may have a hard time acknowledging the health and survivability of their chosen coping mechanisms, we see that the risks continue.

I think that maybe some people around me had a hard time accepting this idea.  Their resistance to this acceptance shaped an important situation that brought me one step closer to studying information security.  Specifically, it shaped the reaction to attacks that I saw one year.  It was the year I found a RAT:  a remote access Trojan, buried deep in a production computer.

\section{First Light}

I was in the back row of a boring sales meeting when I first got the news of a devastating hack that hit one of my clients.  Because of the way some people reacted and some of the things that happened later, I won't tell you who they were.  While not everything turned out for the best for everyone, I don't find it hard to know that there was plenty of good in even the worst people involved.  With that, let's continue.

The tables had those white nylon table cloths.  I had given up my nametag holder so that someone else, more important in the meeting, could have a replacement handy.  While I wanted to be helpful and on hand, I didn't really have a stake in the main mission for these people in this huge sales meeting.  I was the programmer.  As a wizard of the mysterious, I knew just enough of the black arts to be dangerous.  I knew enough of the tediousness of building things right to take care of those I built for.  I knew I would die slowly of dehydration in this sales meeting if it were not for that purloined can of Sprite.

I got a text from my boss who sat on the other side of the room.  It said that we'd been hacked.  Right as the meeting closed, and people milled around in the intermission between speeches, he called me on the phone from the front of the room.  The entire site was wiped out, he said.  Get back to the office.  I was already on my way out the door.

When I got a look at everything that was missing from the directory, I was surprised that some files remained.  A huge swath had been cut into the site.  There were some really big losses.  Where I expected to see files and directories, I saw emptiness.  That evening, I began a methodical review of what had been destroyed and what remained.  I used a guiding scrap of untouched information to remind me what should have been where.  Hours later, we got \gls{backup}s installed.  The site was basically repaired.  We had no idea, really, of what happened to us or why.

After this event was over, in the weeks that followed, we would kick around our guesses.  We never really did have any good idea of exactly what had occurred.  We speculated about motive.  We had reported it to the law, and we had seen them do nothing.  We resolved to build some parts back better, and we did.

About six weeks later, we found the first bad directory.  In a folder of computer files, we found programs that shouldn't be there.  There would be more.  Usually, they stood out to us because they were international; these files were programs that contained code for sending the same message in a variety of languages, based on the receiver's identified preferences.  We didn't do anything international.  Our principle users were confined to domestic business.  Let's just say that they didn't speak French, or Italian, or Danish, or German, or any of those other languages symbolized by those little flag icons we saw in the bad directory.  The \gls{malware} that we found on this hacked computer catered to more languages and cultures than our business normally ever would.

These bad directories, containing collections of flag icons, were templated carding sites.  From reading the code, it was obvious that the purpose of these sites was the deception of people and the capture of their financial information.  We found a new one about every six weeks.  We would look at it carefully, and then rip it out.  A new one would crop up in a different place, some weeks later.  The sites were basically the same, but they would often be perfectly decorated to look like the real thing.  They would look like an actual banking site or a famous login site or some type of digital shopping cart.  Large graphics would fill the backgrounds.  Often, they would be picture-perfect replicas of the real websites users would expect.  A Javascript program, usually a SPA, single page application, would run the forms.  This means that the user's own computer, not some distant machine, would do the computational work of capturing the victim's financial data and send it to the hackers who had set up the carding site.  A few of the carding sites were so sophisticated that they would simulate page clicks; the duped user would probably think that they were visiting a multi-page site, like the one they expected.  Meanwhile, they would stay on the same page all along.  They would dutifully fill out those forms, entering their personal financial information.  It'd never go where they thought.  Instead, it'd be all emailed back to the attacker who ran the carding site.

I learned a lot from reading those pages.  The attackers had installed a fairly good set of Javascript functions to react to different kinds of credit cards.  One type of function would use some simple math to examine the credit card member that had been entered by the victim-user.  Then, the function would classify which kind of credit card had been used, based on the characteristics of the provided numbers.  Using this logic, the site would react to show the user what kind of credit card had been entered.  Their sites could imitate the reactions on commercial websites; type in the card number, and the program would identify what kind of card you were inputting.  Maybe users would be reassured to see the site reacting to their input in a professional way.

Another time, somewhat near the end, I could tell that the site was related to a person who didn't know very well what they were doing.  It flat out looked to me like the program they installed was a lesser version of the others.  It seemed as if that particular carding operator didn't have the technical skill to flesh out and code a good, working design.  We deleted that one, too.

Now, all the while this was going on, for months, we were stifled by our inability to gain root access to control this box.  We couldn't take the normal actions website administrators undertake to monitor, to limit, or to exert power over user actions on the website.  We weren't stymied by criminals.  The previous contractor who had set up the computer had never fully relinquished basic information needed to control the account on this rented server.  In the end, he did give it up.  That was over a year after our first big \gls{attack}.

It took me a long time, months, to find where the server side logs were kept.  These logs were text files automatically written by the computer.  Each time it would receive a call for a file or web page, it would note information about that call.  They described the specific computer address of where every call to the computer had come from.  It showed what file had been accessed each time.  Often, these log entries also detailed commands used by attackers to try to force bad input into the database in order to overwhelm it and grant them an opportunity to cram their own commands into the machine.  These logs were automatically compressed into zipped files.  The hosting company's control panel for the site provided us with a small window on the last 300 entries.  The site was passing so much traffic that 300 log lines was a ridiculously small slice of what was needed.  The account on that leased server was set up to run web pages that received very little traffic, like an amateur's blog.  On those sites, perhaps 300 lines would reveal most of a day's visitors.  For these directories, 300 calls could be made to the site in minutes.  Our client had grown much bigger than its earlier allocation of resources.

When I finally found the logs and unpacked them, a new world awaited.  After some griping, my boss grudgingly admitted that log reviews were helping to protect our website.  I wrote scripts to scrape text copies of the log files.  These programs would automatically read copies of the logs I had downloaded.  Many of the logs were tens of thousands of lines long.  In them, there would be an uncountable number of legitimate calls to the site.  Buried in that collection would be the commands to illegitimate carding sites on the computer.  Also in there would be records of other \gls{attack} attempts.  I soon learned to look for specific keywords.  I saw the common shapes of SQL injection attacks that tried to force commands into the database engine to get it to reveal its contents.  I saw indicators of XSS attacks from referent sites.  Victim visitors who had arrived at the carding sites seemed to have been already compromised by a visit to a previous website.  These calls came in with encoded data already appended to the URL in a way that suggested their requests had been partially processed by another system.  Perhaps their searches on a famous search engine in Russia had indicated they had interests that made them likely victims for these scams.  I picked up IP addresses and ran them through Shodan, a special search engine for the Internet of Things.  Shodan offered the advantage of telling us, by IP address, where in the world a computer was.  Sometimes it would also tell us if other computers were part of a network that it ran.  We could see some information about the types of programs that were on those computers.  It might also include registration information of people and businesses related to the computers.  Sometime there were sample code replies to common types of calls made to those computers.  Cataloging and correlating information from the logs and Shodan proved to be very valuable, over time, to understand the international nature of different kinds of \gls{attack} attempts on the site.  I learned about my attackers.  I could often tell the difference between different styles of attackers.  I could see the plain difference between the lone Metasploit user who was running sqlmap from Tehran and the multi-country curl-based recon attacks from South America and the other calls that were all part of the same Bit Torrent botnets out of Russia and Ukraine.

One thing I didn't see, though, was my main, lone attacker.  He didn't use web addresses to get in.  Somewhere, buried in the directories was a shell.  I found it about a year after he first attacked.

\section{Attack Topology}
This is what I believe he did with that shell.

On arrival, he unpacked his calling card.  This flash, or identifying graphic, was his.  It was basically a large, rectangular picture of a music playing program.  His hacker name was literally written right across it in big, bold letters.  Next, he unpacked a template for the card site he would use over and over again.  All of this was stored in the same directory that held the Trojan program.  This encoded source code could be read by the computer, but was difficult for people to read.  It provided the attacker with full access to read and write files on the compromised computer.  It was his main mechanism for putting his programs on the machine.  It would prove difficult to spot, since it was buried in a large, cluttered directory filled with so many files that it was not practical to inspect them by sight.  With no practical monitoring or real security measures, our attacker was able to conceal his Trojan in a pile of unsorted trash for months.  Our failures were his camouflage.

I don't think we were the only site he had.  My suspicion was that an effective carder would be like an effective check kiter:  he would operate in many locations continuously.  Check kiting, more popular during the time when paper checks were dominant, was a scam that worked like this.  The kiter would go to a bank and write a bad check.  While the check was in the process of clearing, he would have a credit.  The kiter would go to another bank and write another bad check; this time the check would be written against the account where the bad check had been deposited but not yet collected against.  An effective check kiter would have many such transactions in process.  These checks would be "up in the air," figuratively flying like kites.  While I never saw any other the other compromised sites or transactions to them, I felt that it was most likely that websites compromised for carding would be set up like beads on a necklace.  Our carder would visit only every so often; he would set up a new templated site about every month to six weeks; this was not because his business was idle.  Instead, I suspected that we were just one node in a very long line.  By extension, I suspected that the other sites used in other phases of the operation were, too, used like a string of pearls or like beads on a necklace.  One by one, traffic would count their way down the line.  They would control which of the sites in the lines of prepared computers would be used at a given time.  Only a small segment of the compromised machines would need to be exposed at a time.  By using many sites like this, perhaps an attacker could reduce repetition; during my time in the military we were often told the old maxim, "Repetition kills."  Repeated behaviors yield more readily observable patterns.  Our attacker the carder had a lot to hide from.

To help his hiding, he would install blacklists.  These programs, .htaccess Perl scripts, would deny the directory entry, discovery, or communication to sites listed.  Normally, these small programs would be installed in a directory by system administrators to control activity in the directory.  Since he was ensconced on the computer, he provided his own protective site administration.  Later, when I found several of these blacklists, I noticed that my attacker had diligently commented his purpose for blacklisting many of the sites.  In blacklisting, the carding web programs that were in the templated sites would use two layers of defense:  iterative, programmatic denial, and explicit denials.  For instance, the attacker would deny entire ranges of IP addresses.  Large, well-funded multinational businesses would often have big blocks of IP addresses for their use, taking up a whole range of numbers.  Since these companies often had well-staffed and adept computer security departments, he made sure inquiries from their computers would be denied.  Sometimes he would have programs do an iterative comparison of an incoming call's IP address; then he would deny.  That is, he would use his program to receive an address value, then he would loop over a short list of known values that described different limits of ranges to avoid.  This type of \gls{attack} he used often to conceal his operation from larger companies, like Google or Cloudflare, who had well-funded, in-depth defenses.  In big corporations, the defensive security teams would be more likely to be supporting companies that would take significant legal and digital action against the carder.  Also listed, often explicitly, were individual IP addresses of people or other criminal organizations that had evidently been seen as a \gls{threat} to the carding operation.  In these instances, his programs would go down a long list in a simple sequence and compare the caller's IP with known IPs to avoid.  The attacker was literally trying to avoid the competition.  Sometimes the comments next to the IP address would be a rude insult about someone's personality or a note that a successful operation was running at that address.  In all, the blacklist would help the carding operation remain undercover.

I also never saw any indication that the renters, the skiddies, knew where the site was, either.  These skiddies, or script kiddies, were the type of attackers who participated in computer attacks, but did not directly possess enough technical knowledge on their own to conduct an \gls{attack} manually.  Instead, they would buy programs put together by others to conduct an \gls{attack} at the touch of a button.  Perhaps the attacker's very customers were somehow on the blacklist.  I often assumed that the people using the sites that were produced by the template would probably never exactly know where the carding operation was.  While I had no way of knowing, we saw evidence as we went along that maintaining control over information and access was a key part of the carding operation.  The attacker's customers, for example, probably went directly to their own subdirectory.  I just don't believe that they ever used the same back door that the attacker did.  If they had, they would have stumbled upon the templated site.  Then they, too, would be able to replicate it and use it as they will.  Considering that they were probably paying fees for the collected information, a leak like that might have been catastrophic to my attacker's operation.  By not revealing his own method of access and site construction, the attacker was probably employing his own risk controls for his own ends.

The attacker would need someone to rent the site.  These skiddies were probably unable, or unwilling, to run their own carding operation.  I think my attacker had both kinds.  On one occasion, we saw poorly constructed and edited templates; those may have been the less able.  I think the remainder were the unwilling.  I suspect that some of them were unwilling to run their own carding operation because they were running a much larger and involved operation:  it was clear that there was a lot of sexual content to the bait sites that brought some of the people to the carding sites.  Perhaps people in the prostitution business needed operational separation for their own reasons.  Or, maybe, they were low on the technical skills needed to recon, vet, establish, run, conceal and maintain a long string of carding sites on the web as a professional-grade service.  For whatever reason, it was clear that my attacker who unpacked carding sites was also linked up with others who seemed to be his customers.

These skiddies, too, would have baiting sites that also seemed to be like beads in a necklace.  They wouldn't perfectly repeat.  Each day, there would be a slightly different line.  But often, the line of sites would be similar enough to be recognizable.  One set of lines seemed to be from a Russian search engine that was somehow laced with Javascript in a way that implied some cross-site scripting was at work.  Most would be plain calls to carding URLs.  Regardless of structure or sophistication, it was clear that the baiting sites were like one long string of compromised computers, too.  Perhaps they were sometimes legitimate servers that had been directed to meet the carding operations.

These skiddie customers revealed their baiting sites in the server logs.  Over the year that I was learning about these attacks, in the months when I could scrape the logs, I gathered a list of keywords associated with referring sites.  Phrases like "xxx" or ".ru" and many others:  they would be in log entries going to directories that I knew should not normally be receiving traffic.  In turn, I would look up many of the IP addresses associated with those keywords by using Shodan.  I would build spreadsheets filled with IP addresses, calls to URLs, and other information.  I would learn which IPs were associated with locations in foreign countries.  I would learn which IPs were in the same Bit Torrent botnets.  I learned other technical characteristics about various groupings of the websites.  Over time, it was clear that the flood of traffic discovered on some days was aimed directly at a new installation of one of these carding sites.  The logs would show me, by their directories, where to go.  When I got there, sure enough, I would see the familiar flags.  I found slightly modified carding sites each time.  I would study the supporting graphics.  I would read over the associated programs.  I could see that each time, the installation was slightly different.  One social role would be a leader among the referring sites:  sex.

Sexual topics would often be built into the domain names of the baiting sites that referred traffic to the carding sites.  I call them baiting sites because I believe that few people would have believed that they were going to be duped into using these carding sites.  Instead, it seemed that the visitor, the victim's motivation, was sexual activity.

No matter what kind of sex you might desire, these sites might have catered to you.  On one set, the site would be a menu of scantily clad prostitutes.  Women of all kinds would be there.  Some were obviously young and beautiful.  They would wear different kinds of lingerie.  Some would be topless.  Some would not face the camera.  Each woman on a page would have a small block for an ad.  Next to her photo, there would be a price list, often in Russian.  Once I translated these ads using Google translate.  I calculated the exchange rate of the prices.  It was clear that the last three entries on the list of menus were for:  one hour, two hours, or all night.  The all night deluxe experience, the most expensive on the menu, had a going rate equivalent to \$30 \gls{us}.  At the bottom of each ad was the suggestion that these women accepted credit card payments.  The payment methods they accepted would be just like the ones in the carding sites.

Let's not forget that these sites offered many kinds of attractions.  Did you want to start a video chat to talk to two young, suntanned Brazilian men?  Did you want to talk to two girls instead?  Did you favor flirting with the old or young?  Or, were you just shopping?  Were you a pregnant woman, looking to buy clothes for your baby?  Did you want to buy or price an expensive car by submitting your information for a quote?  Did you use online products that required a famous sign-in procedure?  Any of these ploys, not just plain prostitution, could be the bait that got someone referred to carding sites such as these.  Over and over, we saw many different kinds of referent sites in the logs.

So, somehow the skiddie would have made a deal with the attacker to obtain the credit card information of the victim.  The skiddie did not immediately get the information.  Instead, the programs which collected the card information would email them to a public email account.  There, the attacker would probably log in, collect the data, and choose to forward it to his renter-skiddie when he was ready.  The address on the email account matched the flash on the calling card and later claims of attribution.

\section{The Reveal}

When I found the shell, I used some of my contacts to smoke out my guy.  I got independent confirmation on what type of Trojan it was.  Then came the better intel.  Based out of West Africa, my guy had called in, probably through Paris, to \gls{attack} the site.  The calling card he unpacked in the early stages confirmed who he was.  He unfurled a political banner on the site for some defacement.  He bragged about his deeds to some friends.  Then he got to work on building his empire.

During the year I had repeatedly seen some calls, encoded as chars, or single character variable values, that I thought might be the signature of my attacker.  I looked up this code.  I had found some gaming sites, some Facebook pages, some clues in forums about his interests.  At about the right age, in about what I thought was the right place, I thought for sure this would be the calling card of my guy.  I was wrong.  Perhaps those char-encoded values were the signature of another attacker, one who did less damage.  Or maybe it was a magic value or command of some type that would be used by programs buried in the compromised machine.  Whatever it was, I have to admit that one particular call did not seem to match the public face of the person behind the carding operation.  Perhaps our honey attracted a lot of flies.

When my contacts revealed to me the hacker name of the person who attacked the site, I looked him up.  He had even published videos on YouTube showing attacks like those he probably carried out against us.  I remember watching the video carefully, noting the name of the automated tools he used to carry out these attacks.  I made sure I got a copy.

These attacks, and those \gls{attack} tools, became the focus of my research into injection attacks.  I have become curious about evaluating the effectiveness of these automated tools.  Seeing so many attacks, it became clear that some were better than others.  The better the automated tool, I suggested, perhaps the lesser the skill an attacker needed.  While a true blackhat might not need much of anything to carry out his own custom-built attacks, many more uneducated, untrained, and inexperienced skiddies could carry out attacks with these automated tools.  Type in a website address of a target destination and click a button.  For some automated \gls{attack} tools, that's all there was to it.  I have decided to study those.

\markchapterendnotes